% paper/sections/12e_interface_field.tex
%
% Tier IV: 圧力ソルバの検証
%   U6 = Split-PPE + DC + HFE 密度比独立性
%        (a) DC k=2,3 反復収束率
%        (b) Split-PPE 解析 jump-MMS ρ-ratio 1:1 → 1:1000
%        (c) HFE 1D step field extension O(h^6)
% ═══════════════════════════════════════════════════════════════════════════════
\subsection{Tier IV：圧力ソルバ}
\label{sec:verify_solver}

Tier I の CCD-Poisson + DC（U2）と Tier III の界面処理（U4/U5）を組み合わせて
構成される圧力ソルバの中核は，分相 PPE 解法（§~\ref{sec:split_ppe_framework}）
＋ DC 反復（§~\ref{sec:defect_correction_main}）＋ HFE 場延長
（§~\ref{sec:field_extension}）の 3 段構成である．
本 Tier の単体テスト U6 はこの 3 primitive を密度比独立性の観点から
suite 検証する．

% ─────────────────────────────────────────────────────────────────────────────
\subsubsection{U6：Split-PPE + DC + HFE 密度比独立性}
\label{sec:U6_split_ppe_dc_hfe}
% ─────────────────────────────────────────────────────────────────────────────

\noindent\textbf{目的：}
分相 PPE + DC $k = 3$ + HFE 構成が，密度比 $\rho_l/\rho_g = 1$--$1000$ に
わたって等密度時の高次精度（$\Ord{h^7}$ Dirichlet 超収束 / $\Ord{h^5}$ Neumann 律速）を
維持することを 3 sub-test で実証する．

\noindent\textbf{Part 2 対応 primitive：}
\#22 (Split-PPE),\#23 (HFE),\#24 (DC iteration)．

\noindent\textbf{下流連関：}
DC convergence vs $\rho$-ratio（§~\ref{sec:varrho_dc_convergence}），
高密度比 droplet．

% ───────────────
\paragraph{(a) DC $k = 2, 3$ 反復収束率}
\label{sec:U6_dc_iteration_rate}

§~\ref{sec:defect_correction_main} の DC 反復が密度比に応じてどの程度
収束率を保つかを，$\omega$ 緩和との結合で診断する．
等密度（$\rho_l/\rho_g = 1$）では $\omega \in [0.2, 0.7]$ で収束するが，
$\rho_l/\rho_g \geq 5$ では一括 PPE のすべての $\omega$ で停滞することが
旧解析（§~\ref{sec:dc_iteration_accuracy}（a））で判明している．
分相 PPE では各相内が等密度となるため，$k = 3$ で $\Ord{h^7}$ Dirichlet 超収束
（液相内部）が回復する（次の sub-test (b)）．

\noindent\textbf{結論：}DC は分相 PPE と組み合わせることで密度比に独立な収束率を達成．\textbf{合格}．

% ───────────────
\paragraph{(b) Split-PPE 解析 jump-MMS：$\rho$-ratio 1:1 → 1:1000}
\label{sec:verify_split_ppe_dc}
\label{sec:split_ppe_recovery}%   backward compat
\label{sec:gfm_ppe_recovery}%     backward compat

製造解 $p^* = \sin\pi x\,\sin\pi y$（Dirichlet BC），
円形界面（$R = 0.25$，中心 $(0.5, 0.5)$），
smoothed Heaviside 密度場（$\varepsilon = 1.5h$）に対して
3 手法を比較する：
(a) 変密度一括 DC $k = 3$，
(b) 分相 PPE + FD 直接法（$\Ord{h^2}$），
(c) 分相 PPE + DC $k = 3$．
液相内部（$\phi > 3h$）で $L^\infty$ 誤差を計測する．

\begin{table}[htbp]
\centering
\caption{3 手法の PPE 精度比較（$N = 64$，液相内部 $L^\infty$）}
\label{tab:split_ppe_dc_sweep}
\renewcommand{\arraystretch}{1.2}
\footnotesize
\begin{tabular}{rccc}
\toprule
$\rho_l/\rho_g$ & 変密度一括 DC & 分相 PPE + FD & 分相 PPE + DC $k{=}3$ \\
\midrule
   1 & $6.6 \times 10^{-12}$  & $2.0 \times 10^{-4}$ & $9.4 \times 10^{-11}$ \\
   2 & $8.4 \times 10^{-8}$   & $2.0 \times 10^{-4}$ & $9.4 \times 10^{-11}$ \\
   5 & $5.8 \times 10^{-7}$   & $2.0 \times 10^{-4}$ & $9.4 \times 10^{-11}$ \\
  10 & $2.9 \times 10^{-7}$   & $2.0 \times 10^{-4}$ & $9.4 \times 10^{-11}$ \\
 100 & $9.1 \times 10^{-6}$   & $2.0 \times 10^{-4}$ & $9.4 \times 10^{-11}$ \\
1000 & $2.7 \times 10^{-5}$   & $2.0 \times 10^{-4}$ & $9.4 \times 10^{-11}$ \\
\bottomrule
\end{tabular}
\end{table}

\begin{table}[htbp]
\centering
\caption{格子収束（$\rho_l/\rho_g = 1000$，液相内部 $L^\infty$）}
\label{tab:split_ppe_dc_convergence}
\renewcommand{\arraystretch}{1.2}
\footnotesize
\begin{tabular}{rcccccc}
\toprule
$N$ & 一括 DC & 次数 & FD 直接 & 次数 & DC $k{=}3$ & 次数 \\
\midrule
  8 & $2.3 \times 10^{-3}$  & ---      & $1.3 \times 10^{-2}$ & ---   & $2.1 \times 10^{-4}$  & --- \\
 16 & $1.9 \times 10^{-4}$  & $3.6$    & $3.2 \times 10^{-3}$ & $2.0$ & $1.5 \times 10^{-6}$  & $7.1$ \\
 32 & $1.5 \times 10^{-5}$  & $3.7$    & $8.0 \times 10^{-4}$ & $2.0$ & $1.2 \times 10^{-8}$  & $7.0$ \\
 64 & $2.7 \times 10^{-5}$  & $-0.8$   & $2.0 \times 10^{-4}$ & $2.0$ & $9.4 \times 10^{-11}$ & $7.0$ \\
128 & $4.1 \times 10^{-5}$  & $-0.6$   & $5.0 \times 10^{-5}$ & $2.0$ & $7.2 \times 10^{-13}$ & $7.0$ \\
\bottomrule
\end{tabular}
\end{table}

分相 PPE + DC $k = 3$ の液相内部誤差は $9.4 \times 10^{-11}$（$N = 64$）で
\textbf{全密度比に対し同一}．
$\rho_l/\rho_g = 1$（等密度，§~\ref{sec:verify_ppe} の $1.29 \times 10^{-10}$）と同水準．
変密度一括 DC は $N \geq 64$ で誤差が増大（負の収束次数；
界面幅 $\varepsilon = 1.5h$ の縮小に伴う $1/\rho$ ジャンプの急峻化が原因）．
$N = 128, \rho = 1000$ では分相 PPE + DC $k=3$ が一括 DC 比 $5.7 \times 10^7$ 倍の精度．

\begin{tcolorbox}[resultbox, title={分相 PPE + DC $k = 3$ の精度検証}]
\begin{enumerate}[noitemsep]
  \item 液相内部で全密度比（$\rho_l/\rho_g = 1$--$1000$）に対し $\Ord{h^{7.0}}$（Dirichlet）
  \item 相内部の格子収束次数は密度比に\textbf{非依存}（$N = 64 \to 128$ で次数 $7.04$）
  \item 変密度一括 DC 比で $N = 128, \rho = 1000$ にて $5.7 \times 10^7$ 倍改善
\end{enumerate}
\end{tcolorbox}

\noindent\textbf{結論：}分相 PPE + DC $k = 3$ は密度比に非依存な $\Ord{h^{7.0}}$（Dirichlet）/ $\Ord{h^{5.0}}$（Neumann BC 律速；§~\ref{sec:verify_ppe_neumann}）．\textbf{合格}．

% ───────────────
\paragraph{(c) HFE 1D step field extension $\Ord{h^6}$}
\label{sec:verify_field_extension}
\label{sec:verify_gfm}%   backward compat

§~\ref{sec:field_extension} の HFE が，
ジャンプ補正済み片側データを界面越しに $\Ord{h^6}$ で延長することを
1D 製造解で検証する（風上差分経路は密度ジャンプ越え NaN により廃止済み）．

\noindent\textbf{結果：}
1D Hermite テストでは $N = 32 \to 128$ で
$L^\infty$ 誤差が $9.76 \times 10^{-10} \to 5.97 \times 10^{-14}$ に減少し，
実測次数 $p \approx 7.0$ を示した（$N = 256$ で丸め誤差床）．
2D テンソル積逐次補間は $N = 32 \to 128$ で $p = 6.33, 5.79$ の高次収束を示すが，
$N = 256$ では $L^\infty = 4.59 \times 10^{-11}$ で丸め・補間誤差床（$p = 2.69$）に到達．

\noindent\textbf{結論：}1D で $\Ord{h^{7.0}}$，2D で $N \leq 128$ まで $\Ord{h^{5.8}}$ 以上．
\textbf{合格}（1D）／ $\triangle$（2D は丸め床到達）．
CSF 正則化 $\Ord{h^2}$ より十分小さい誤差水準には到達しており下流非律速．
